\bookchapter{Глава 32: Экстренная помощь по безопасности телефона}

\chapterintro{Эта глава — «план на случай чрезвычайной ситуации». Надеемся, он вам не понадобится, но если телефон потеряется, украдут или обманут, памятка поможет минимизировать убытки. При землетрясениях, ливнях и других бедствиях телефон тоже остаётся инструментом для приёма предупреждений, вызова помощи и связи. Прочитайте сейчас, сохраните скриншоты ключевых шагов или распечатайте и держите дома: в критический момент телефона может не оказаться под рукой, а во время стихии он может быть разряжен или без сети.}

\booksection{Сценарий 0: землетрясение / экстремальная погода — как телефон помогает выжить}

Речь о безопасности людей, а не денег на счетах. Настройки предупреждений — в главе 13 «Погодные, стихийные и сейсмические предупреждения»; здесь — что делать после сигнала.

\textbf{Сигнал землетрясения: сначала укрыться, потом эвакуация — не снимайте видео}\par

\begin{enumerate}
  \item Укройтесь рядом: упасть, укрыться, удержаться (под столом, у несущей стены, у прочной мебели), защитите голову и шею; держитесь подальше от стёкол, внешних стен и люстр.
  \item Не бегите слепо по лестнице: основной толчок длится от нескольких секунд до десятков, лестничные клетки могут быть опаснее; лифт категорически нельзя.
  \item После толчка эвакуируйтесь организованно: по лестнице, следите за афтершоками и падающими предметами; выходите на открытое место, подальше от столбов, щитов и аварийных зданий.
  \item Телефон — по делу: короткое сообщение экстренным контактам («Я на XX, в безопасности/ранен, нужна помощь») экономит батарею лучше длинных звонков; если есть SMS/WeChat — пишите текстом.
  \item Нет сигнала: смотрите главу 7 «Спутниковый SOS / спутниковые сообщения» (если устройство поддерживает); или попробуйте на возвышенности и открытом месте; в некоторых случаях помогает радиосвязь Bluetooth без сети.
  \item Отключите лишнее: включите авиарежим, затем по таймеру включайте/выключайте, выключите сканирование Bluetooth, уменьшите яркость — так продержитесь дольше; последний заряд оставьте на сообщения о безопасности и спасении.
\end{enumerate}

\begin{tcbwarn}{Ливень / тайфун / селевой поток: действуйте уже на этапе предупреждения}
\begin{itemize}
  \item Увидели красное предупреждение: по возможности не выходите, держитесь подальше от подземных пространств, путепроводов, русел и крутых склонов.
  \item Завалило на улице: поднимайтесь в прочные здания, не вброд через воду (в ней могут быть открытые люки и утечки тока); звоня в 119/110, назовите ориентиры и уровень воды.
  \item Перед поездкой: проверьте риск в погодном приложении и местных каналах предупреждений; на дальних автопутешествиях избегайте горных дорог у воды и обрывов.
  \item Семейный протокол: договоритесь «если связь пропала, ориентируемся на последнее место по WeChat/SMS», чтобы не разряжать батарею бесконечными звонками.
\end{itemize}

\end{tcbwarn}

\begin{tcbpractice}{4 вещи, которые стоит сделать заранее}
\begin{enumerate}
  \item Включите системное предупреждение о землетрясении / надёжное приложение и разрешите ему «прорываться» сквозь беззвучный режим.
  \item Закрепите в контактах 2–3 «экстренных контакта», они могут отображаться на экране блокировки (медицинская карта iOS/Android или информация на экране блокировки).
  \item Держите повербанк дома и в сумке; перед дальней поездкой зарядите его и возьмите кабель.
  \item Запомните хотя бы один номер, по которому можно позвонить без контактов (мобильный родственника), запишите на бумагу и держите в кошельке.
\end{enumerate}

\end{tcbpractice}

\booksection{Сценарий 1: телефон потерян/украден — золотые 30 минут}

Первые 30 минут после пропажи — «золотое время». Правильные действия помогут сохранить деньги и приватность. Действуйте по порядку, не пропускайте шаги:

\textbf{Стандартная процедура первых 30 минут}\par

\begin{enumerate}
  \item Позвоните на свой номер с чужого телефона (1-я минута): возможно, он просто упал рядом, и добросовестный человек поднимет.
  \item Удалённо заблокируйте телефон (2–5-я минута): с чужого телефона или ПК зайдите в сервис поиска: iPhone: iCloud.com > Найти > выберите устройство > Отметить как утерянное
  \item Android по брендам: Huawei cloud.huawei.com / Xiaomi i.mi.com / OPPO cloud.oppo.com и др.
  \item Заблокируйте SIM (5–10-я минута): позвоните оператору (China Mobile 10086 / China Unicom 10010 / China Telecom 10000) и попросите приостановить обслуживание SIM. Это ключевой шаг — с вашей SIM мошенник получает коды подтверждения и может сбросить все пароли.
  \item Заморозьте платёжные аккаунты (10–20-я минута): Alipay: 95188 — блокировка, или на чужом телефоне войдите в Alipay > Мой > Настройки > Центр безопасности > Экстренная заморозка
  \item WeChat: после вызова 110 свяжитесь с поддержкой WeChat или в Центре безопасности WeChat > Заморозить аккаунт
  \item Банковское приложение: позвоните по номеру на обороте карты и попросите заморозить интернет-банк
  \item Предупредите близких (20–25-я минута): напишите в Moments или попросите родных сообщить о пропаже, чтобы мошенники не обманули других от вашего имени.
  \item Заявление в полицию (25–30-я минута): позвоните в 110 или обратитесь в ближайший участок и получите талон о принятии заявления. Он понадобится для страховки и споров с банком.
\end{enumerate}

\booksection{Сценарий 2: вас обманули — 72 золотых часа}

Если вы поняли, что вас обманули, главное — сохранять спокойствие. Паника ведёт ко вторичным потерям. Действуйте так:

\begin{enumerate}
  \item Немедленно прекратите все переводы: что бы мошенник ни говорил, не переводите ни копейки. «Ещё один залог — и вернём предыдущие деньги» — классический приём повторного обмана.
  \item Сохраните все улики: не удаляйте переписку и журнал звонков! Сделайте скриншоты переписки, переводов и данных счёта собеседника. Это ключевые доказательства для заявления и возврата денег.
  \item Немедленно позвоните в 96110: это общегородская линия борьбы с мошенничеством. Оператор подскажет экстренную остановку платежа — если перевод недавний (обычно в течение 24 часов), банк может заморозить счёт получателя.
  \item Позвоните в 110: возьмите телефон, паспорт и платёжные документы, подайте заявление в ближайший участок.
  \item Свяжитесь с банком/платёжной платформой: сообщите об обмане и попросите заморозить счета по делу. Если с кредитки украли деньги, попросите банк запустить процедуру chargeback.
  \item Смените все пароли: на другом безопасном устройстве смените пароли Apple ID/аккаунта бренда, WeChat, Alipay и банковского приложения и включите двухфакторную защиту.
\end{enumerate}

\booksection{Сценарий 3: телефон заражён вирусом/трояном — что делать?}

Если появились странные симптомы — непонятная реклама, автоматическая установка приложений, резкий рост счетов за связь, аномальный расход трафика — вероятно, телефон заражён:

\begin{itemize}
  \item Шаг 1: войдите в безопасный режим — на Android удерживайте кнопку питания для выключения, снова включите и после вибрации удерживайте клавишу уменьшения громкости, пока не появится «Безопасный режим» (в левом нижнем углу). В безопасном режиме сторонние приложения отключены — удобно искать причину.
  \item Шаг 2: удалите подозрительные приложения — в безопасном режиме откройте Настройки > Приложения, отсортируйте по дате и удалите недавно установленные незнакомые.
  \item Шаг 3: установите антивирус и просканируйте — скачайте «Tencent Mobile Manager» или «360 Mobile Guard» и сделайте полное сканирование.
  \item Шаг 4: сброс до заводских настроек (крайняя мера) — если ничего не помогло, сохраните важные данные и сделайте сброс. Внимание: сброс удалит всё, сначала обязательно сохраните фото, контакты и другие файлы.
  \item iPhone почти не заражается — песочница iOS делает выживание вредоносов крайне сложным. Если iPhone ведёт себя странно, скорее всего, это баг приложения или настройки, а не вирус. Сброс до заводских решает 99\% проблем.
\end{itemize}

\booksection{Сценарий 4: аккаунт взломан — быстро верните контроль}

\begin{enumerate}
  \item Немедленно смените пароль: если ещё можете войти, смените пароль сразу. Новый пароль должен полностью отличаться от старого.
  \item Отключите чужие устройства: в настройках безопасности найдите «Управление устройствами входа» и отключите все незнакомые.
  \item Проверьте журнал входов: за последние 7 дней, если есть входы с других мест — аккаунт действительно взломан.
  \item Включите двухфакторную защиту: если ещё не включена — включите прямо сейчас. Это самый эффективный способ не потерять аккаунт снова.
  \item Проверьте привязки: не изменились ли телефон и почта. Если изменились — восстановите аккаунт через апелляцию.
  \item Проверьте платежи: если взломали WeChat/Alipay, посмотрите, не изменены ли автоплатежи и платежи без пароля, не добавлены ли карты в чужие аккаунты.
\end{enumerate}

\booksection{Сценарий 5: пришло шантажное сообщение — не бойтесь и не платите}

Два типичных сценария шантажа: (1) «У меня твои компрометирующие фото/видео, не заплатишь — разошлю всем в контактах»; (2) «Твой телефон взломан, не заплатишь — заблокирую». Что делать:

\begin{itemize}
  \item 99,9\% — это обман: у них на самом деле нет ваших данных. Рассылая одинаковые угрозы десяткам тысяч людей, всегда найдутся те, кто испугается и заплатит.
  \item Не отвечайте и не платите: ответ означает «этот номер активен», вас пометят как «подходящую цель» и будут донимать дальше.
  \item Сделайте скриншот и сохраните улики, затем заблокируйте собеседника.
  \item Смените пароль: для собственного спокойствия обновите пароли основных аккаунтов.
  \item Если действительно есть риск утечки приватных данных: обращайтесь в полицию. У правоохранителей есть опыт таких дел.
\end{itemize}

\booksection{Сценарий 6: ИИ-подмена лица/голоса — глазам больше верить нельзя}

С 2024 года ИИ-подмена лица и голоса позволяет мошенникам в видеозвонке «стать» вашими родными, начальником или другом. «Мам, у меня беда, срочно переведи деньги» — на видео человек выглядит и звучит как ваш сын, но это подделка. Один из самых опасных и трудноуловимых приёмов.

\textbf{Шаг 1: ищите слабые места — четыре типичных признака подделки}\par

\begin{itemize}
  \item Частота моргания: живой человек моргает 15–20 раз в минуту, ИИ-подмена часто моргает слишком редко или ритм скованный.
  \item Губы и звук не совпадают: между движениями губ и голосом ощущается задержка в доли секунды.
  \item Размытые края лица: контур, область ушей, линия волос могут «рябить» или мерцать.
  \item Несогласованный свет: блики и тени на лице не совпадают с освещением фона, а при повороте профиль «ведёт» черты некорректно.
\end{itemize}

\begin{tcbinfo}{Шаг 2: три способа проверить собеседника}
\begin{itemize}
  \item Попросите сделать действие: повернуться, закрыть лицо рукой, прикрыть рот и сделать движение. При масках и резких поворотах ИИ чаще всего «выдаёт» себя — лицо плывёт или не поспевает.
  \item Задайте вопрос, известный только собеседнику: придумайте ложный вопрос, например «в какую клинику возили нашу собаку в прошлый раз?» — мошенник не ответит или начнёт путаться.
  \item Перезвоните на известный номер: сбросьте видеозвонок и позвоните на номер из контактов. Мошенник не контролирует настоящий номер — обман раскроется сразу.
\end{itemize}

\end{tcbinfo}

\begin{tcbwarn}{Шаг 3: если обман подтверждён — порядок действий}
\begin{enumerate}
  \item Немедленно прекратите переводы: даже если «горит», остановитесь. Те, кто торопит «скорее переводи и не клади трубку», почти всегда мошенники.
  \item Сохраните все улики: запись экрана, скриншоты видеозвонка, переписки, счёта и переводов — это материал для экспертизы.
  \item Позвоните в 96110: сообщите, что это мошенничество с ИИ-подменой лица/голоса, оператор запустит экстренную остановку платежа.
  \item Позвоните в 110: с полным пакетом улик обратитесь в участок, особенно подчеркните использование ИИ-фальсификации — это поможет классифицировать дело.
  \item Предупредите близких: мошенники могут использовать вашу подделку против других — заранее напишите в Moments или семейном чате.
\end{enumerate}

\end{tcbwarn}

\booksection{Сценарий 7: SIM swap (угон SIM) — ваш номер «украли»}

Малоизвестная, но крайне опасная атака. SIM swap — когда злоумышленник подделывает ваши документы и оформляет дубликат SIM в салоне. После этого ваш сигнал пропадает, а на новой карте атакующего приходят все SMS-коды: Alipay, банки, WeChat — все двери открыты.

\begin{tcbinfo}{Как заметить? Три сигнала SIM swap}
\begin{itemize}
  \item Внезапно «Нет сети» / «Нет SIM», хотя вы её не трогали и перезагрузка не помогает — часто момент угона.
  \item Пришло SMS от оператора «Успешная замена SIM» или «Услуга оформлена», хотя вы ничего не делали.
  \item Вдруг перестали приходить коды подтверждения, и другие SMS тоже — потому что их увёл атакующий.
\end{itemize}

\end{tcbinfo}

\begin{tcbwarn}{Защита — «запереть» и SIM-карту}
\begin{itemize}
  \item Задайте PIN-код SIM: путь Настройки > Телефон > PIN-код SIM. После включения при каждой перезагрузке или смене карты нужен PIN. Атакующий не сможет пользоваться картой даже при физическом доступе.
  \item Включите вторичную проверку оператора: в приложении China Mobile/Unicom/Telecom найдите «Центр безопасности» или «Проверка при обслуживании», включите уведомления о чужих заменах SIM и повторный пароль на услуги. Тогда при попытке чужой замены сотруднику потребуется дополнительная проверка личности.
  \item Для важных аккаунтов не полагайтесь только на SMS: банки, Apple ID и другие ценные учётки лучше защищать аутентификатором (Google Authenticator, Microsoft Authenticator) или аппаратным ключом.
  \item Водяные знаки на копиях паспорта: на копиях для оформления пишите «Только для XX» — атакующие часто используют утёкшие паспортные данные для подмены личности.
\end{itemize}

\end{tcbwarn}

\textbf{Действия после угона — скорость решает всё}\par

\begin{enumerate}
  \item Немедленно заблокируйте SIM через оператора: China Mobile 10086 / China Unicom 10010 / China Telecom 10000, попросите приостановить обслуживание и заморозить замену SIM, чтобы остановить атакующего.
  \item Войдите в ключевые аккаунты с других устройств: с ПК или телефона родных немедленно смените пароли банков, Alipay, WeChat, Apple ID и отключите незнакомые устройства.
  \item Заморозьте платёжные аккаунты: Alipay 95188, Центр безопасности WeChat, банковская поддержка — заморозка интернет-банка.
  \item С паспортом в салон: лично в фирменном салоне оформите новую SIM и восстановите номер.
  \item Заявление в полицию: SIM swap квалифицируется как «посягательство на персональные данные граждан». Возьмите улики (время блокировки, убытки, записи о замене SIM).
  \item Полная проверка: после восстановления номера просмотрите недавние входы и финансовые операции на предмет чужого использования.
\end{enumerate}

\begin{tcbwarn}{Профилактика лучше лечения — 5 ежедневных шагов}
\begin{enumerate}
  \item Пароль блокировки экрана + двухфакторная защита — включите оба
  \item Убедитесь, что «Найти мой телефон» включён
  \item Задайте PIN-код SIM (Настройки > Телефон > PIN-код SIM), при потере телефона чужие не смогут пользоваться картой
  \item Не ставьте приложения с неизвестных источников, не открывайте чужие ссылки
  \item Регулярно делайте резервные копии важных данных (фото, контакты, переписки)
\end{enumerate}

\end{tcbwarn}

\begin{tcbpractice}{Практика: сделайте свою «экстренную карточку контактов»}
\begin{enumerate}
  \item Возьмите лист бумаги и запишите: службу оператора (China Mobile 10086 / China Unicom 10010 / China Telecom 10000), банковскую поддержку, линию 96110, полицию 110, телефоны близких.
  \item Запишите почту регистрации Apple ID/аккаунта бренда и резервный контакт.
  \item Держите лист в кошельке или на холодильнике — не только в телефоне, потому что в критический момент его может не быть рядом.
  \item Проверьте прямо сейчас: включён ли «Найти мой телефон» и задан ли PIN-код SIM.
\end{enumerate}

\end{tcbpractice}

\begin{tcbsummary}{Итоги главы}
\begin{itemize}
  \item Золотые 30 минут при потере телефона: позвонить на свой номер удалённо заблокировать заблокировать SIM заморозить платежи предупредить близких заявление в полицию
  \item 72 часа после обмана: прекратить переводы сохранить улики позвонить в 96110 заявление в полицию связаться с банком сменить пароли
  \item Заражение телефона: безопасный режим удалить подозрительные приложения антивирус сброс до заводских (крайняя мера)
  \item Взлом аккаунта: смена пароля отключение устройств проверка входов включение 2FA проверка привязок проверка платежей
  \item Шантажные сообщения на 99,9\% — обман: не отвечайте, не платите, скриншот — и в блокировку
  \item ИИ-подмена лица/голоса: смотрите на моргание, губы, края, свет; просите повернуться/закрыть лицо, задавайте личные вопросы, перезванивайте на известный номер
  \item SIM swap: внезапное «Нет сети» — тревожный сигнал; задайте PIN SIM, включите вторичную проверку оператора, переведите важные аккаунты на аутентификаторы
  \item 5 ежедневных мер: блокировка экрана + 2FA, «Найти мой телефон», PIN SIM, без сомнительных приложений, регулярные резервные копии
\end{itemize}

\end{tcbsummary}
